\section*{Chapitre \ref{ch:ml:online-learning-and-drift} --- Apprentissage en ligne et dérive}

\begin{solution}{exo:ml:online-learning-and-drift:1}
Mémoires $1/(1-\lambda)$~: 100, 500 et 2\,000 observations. Poids relatifs $\lambda^{500}$~: $0.99^{500} = 0.0066$,
$0.998^{500} = 0.37$, $0.9995^{500} = 0.78$.
\end{solution}

\begin{solution}{exo:ml:online-learning-and-drift:2}
$20\times1\,000/252 = 79.4$. Les 74, 78 et 91 du chapitre sont dans le bruit de Poisson de ce comptage ($\pm9$), ce qui est
précisément ce qu'une calibration sur un autre flux sans rupture peut garantir.
\end{solution}

\begin{solution}{exo:ml:online-learning-and-drift:3}
La statistique ajoute $\mu_0 - v - k$ à chaque pas, en moyenne $0 + 0.38 - 0.1 = 0.28$~; partant de zéro, il lui faut $9/0.28\approx32$
pas. La médiane mesurée est de 17 parce que la statistique ne part pas de zéro~: sans rupture, une marche réfléchie avec
dérive $-0.1$ et variance unitaire se situe en moyenne autour de $1/(2\times0.1) = 5$, de sorte qu'à la rupture elle est typiquement déjà de quelques unités
au-dessus, et le bruit avance encore le premier franchissement.
\end{solution}

\begin{solution}{exo:ml:online-learning-and-drift:4}
L'erreur d'un estimateur à oubli comporte deux composantes~: le bruit d'estimation, qui décroît quand la mémoire $M$ croît (environ $p/M$ en
unités de R-carré), et la péremption, qui croît avec la part de la mémoire passée dans des régimes révolus (environ $M/D$ fois le
signal). Minimiser $p/M + cM/D$ donne $M\propto\sqrt{pD/c}$~: la meilleure mémoire croît comme la racine carrée de la longueur du régime
et non proportionnellement à celle-ci. L'optimum du chapitre, 200 observations pour $D = 100$ et 500 pour $D = 2\,000$, s'accorde avec cette
croissance lente~; pour $D = 10\,000$, le flux ne contient que deux régimes et l'absence d'oubli l'emporte.
\end{solution}

\begin{solution}{exo:ml:online-learning-and-drift:5}
Il a été calibré et ses délais mesurés sur des renversements de signe, où les gains tombent vers $-0.38$ écarts-types~; lorsque les
coefficients sont tirés à nouveau au hasard, les gains de l'ancien modèle ne tombent qu'autour de zéro, soit un changement de moitié moindre, et un
seuil de 9 avec une tolérance de 0.1 exige une longue suite de telles observations. Recalibrer la tolérance pour le changement qui compte
(qui s'estompe vers zéro), surveiller une statistique qui y réagit (la corrélation glissante entre prévision et résultat), et garder un
réentraînement calendaire servant de plancher sous le détecteur.
\end{solution}

\begin{solution}{exo:ml:online-learning-and-drift:6}
Le seuil a été réglé sur les deux changements sur lesquels il est ensuite jugé~: son délai est mesuré dans l'échantillon (in-sample) et son
taux de fausses alarmes n'est pas mesuré du tout. Deux événements ne permettent de calibrer quoi que ce soit. Fixer le seuil sur des données sans rupture pour un taux de fausses alarmes
choisi, puis mesurer le délai sur des changements que le seuil n'a jamais vus (injectés, ou tirés d'un historique ultérieur).
\end{solution}

\begin{solution}{exo:ml:online-learning-and-drift:7}
Avec des régimes de 500 pas en moyenne~: jamais réentraîné $-3.68~\%$~; tous les 250 pas $-0.60~\%$ (75 réentraînements)~; avec le \omterm{def:ml:online-learning-and-drift:detector}{test de
Page--Hinkley} $-8.24~\%$ (45 réentraînements)~; \omterm{def:ml:online-learning-and-drift:rls}{moindres carrés récursifs} avec $\lambda = 0.995$ 2.24~\%~; la valeur vraie 5.35~\%. La fenêtre
de 500 observations du calendrier couvre désormais un changement de régime la plupart du temps, si bien que ses ajustements sont périmés~; le détecteur déclenche
beaucoup plus souvent, et chaque fois il réajuste sur les 100 observations écoulées depuis l'alarme, soit cinq coefficients tirés d'un échantillon trop petit pour un signal
valant 5~\% de la variance, et ces ajustements bruités perdent davantage qu'un ajustement périmé. Seul l'estimateur dont la mémoire convient aux
régimes, soit environ 200 observations, reste positif.
\end{solution}

\begin{solution}{exo:ml:online-learning-and-drift:8}
Avec $\lambda = 1$, la mise à jour de $P$ est la formule de Sherman--Morrison pour $P_t^{-1} = P_{t-1}^{-1} + x_tx_t^\top$, donc
$P_t^{-1} = I/\delta + \sum_{s\le t}x_sx_s^\top$. Par récurrence, $P_t^{-1}\hat\beta_t = \sum_{s\le t}x_sy_s$~: la relation vaut à
$t = 0$ avec $\hat\beta_0 = 0$, et $P_t^{-1}\hat\beta_t = P_t^{-1}\hat\beta_{t-1} + P_t^{-1}k_t(y_t - x_t^\top\hat\beta_{t-1})$
avec $P_t^{-1}k_t = x_t$, on obtient $P_{t-1}^{-1}\hat\beta_{t-1} + x_ty_t$. D'où $\hat\beta_t = (X^\top X + I/\delta)^{-1}X^\top y$,
l'estimateur de ridge de pénalité $1/\delta$, qui tend vers les moindres carrés ordinaires quand $\delta\to\infty$.
\end{solution}

\begin{solution}{pb:ml:online-learning-and-drift:1}
\textbf{Partie I.}
\begin{enumerate}
\item 4.9~\% à 5.4~\% selon le flux. Avec des régimes de 100 pas, le modèle doit estimer cinq coefficients à partir de
une centaine d'observations avant qu'ils ne changent~; son bruit d'estimation est aussi grand que le signal.
\item $\lambda = 0.995$ (mémoire 200) pour $D = 100$ et 500, avec un score de 0.75~\% et 2.44~\%~; $\lambda = 0.998$ (500) pour
$D = 2\,000$, 3.71~\%~; $\lambda = 1$ pour $D = 10\,000$, 5.21~\% contre la valeur vraie de 5.29~\%.
\item 0.57~\%, 2.05~\%, 3.69~\% et 4.84~\%. Un \omterm{def:ml:trees-and-boosting:boosting}{taux d'apprentissage} fixe pondère les anciennes observations géométriquement, comme un \omterm{def:ml:online-learning-and-drift:rls}{facteur
d'oubli}, mais le long d'une seule direction de gradient à la fois~; c'est une version moins coûteuse et plus bruitée des \omterm{def:ml:online-learning-and-drift:rls}{moindres carrés récursifs}.
\item L'hypothèse selon laquelle les vrais coefficients suivent une marche aléatoire, de variance par pas fixée par $\lambda$~: choisir $\lambda$ revient à
décider à quelle vitesse on croit que la vraie valeur se déplace.
\end{enumerate}
\textbf{Partie II.}
\begin{enumerate}[resume]
\item Le gain $\hat y_ty_t$, standardisé sur la période d'échauffement. C'est ce que le desk tire du modèle, et sa moyenne change
de signe lorsque la relation s'inverse~; l'erreur quadratique est dominée par le bruit, qui ne change pas.
\item $h = 9$ pour CUSUM et Page--Hinkley (marge et tolérance 0.1), $\delta = 0.0032$ pour ADWIN.
\item CUSUM 20.8 (médiane 17), Page--Hinkley 27.4 (18.5), ADWIN 42.5 (26.5) pas.
\item 74, 78 et 91 sur 20\,000 pas, contre 79 attendus~: conforme à la calibration. Qu'une alarme par an soit trop relève
de la question de coût du desk (partie IV).
\end{enumerate}
\textbf{Partie III.}
\begin{enumerate}[resume]
\item Jamais réentraîné $-7.81~\%$~; tous les 250 pas 2.62~\%~; sur une alarme de Page--Hinkley $-2.67~\%$ (deux réentraînements)~; \omterm{def:ml:online-learning-and-drift:rls}{moindres carrés récursifs} avec $\lambda = 0.998$
3.24~\%~; la valeur vraie 4.62~\%.
\item Ses coefficients appartiennent au premier régime~; dans les autres régimes, ils sont sans rapport avec la vérité, voire opposés à elle, et une
prévision qui se trompe de sens obtient un score inférieur à celui de la prévision nulle.
\item Il a été réglé sur des renversements, alors que le flux retire ses coefficients~; voir l'\cref{exo:ml:online-learning-and-drift:5}.
\item Il réentraîne à chaque pas avec la mémoire qui convient aux régimes, sans seuil à franchir et sans écarter aucun échantillon.
\end{enumerate}
\textbf{Partie IV.}
\begin{enumerate}[resume]
\item \emph{La demi-vie d'un modèle.} Le \omterm{def:ml:online-learning-and-drift:rls}{facteur d'oubli} qui minimise la perte préquentielle croît avec la
longueur du régime~: $\lambda = 0.995$ pour $D = 100$ et 500, 0.998 pour 2\,000, 1 pour 10\,000. À une fausse alarme par 252 pas, le
délai médian avant un renversement de signe est de 17 pas pour CUSUM, 18.5 pour Page--Hinkley et 26.5 pour ADWIN.
\item Un estimateur en ligne lorsque le modèle est linéaire ou a une dernière couche linéaire, avec la mémoire choisie sur l'historique~; un
réentraînement calendaire pour le reste~; un détecteur sur les gains pour déclencher un réentraînement anticipé et une revue humaine.
\item Balayer le \omterm{def:ml:online-learning-and-drift:rls}{facteur d'oubli} ou la fenêtre glissante sur l'historique et lire la longueur du régime sur l'optimum
(voir l'\cref{exo:ml:online-learning-and-drift:4}), ou compter les ruptures structurelles que trouvent les tests du livre~7.
\item Une fausse alarme coûte un réentraînement (calcul, validation, renouvellement du modèle) et la confiance~; un changement manqué coûte les
pertes d'un modèle erroné jusqu'au réentraînement suivant, qui, pour un signal renversé, équivaut à la perte espérée de toute la position.
\item Le livre~7 a mesuré a posteriori la décroissance d'un signal~; le \omterm{def:ml:online-learning-and-drift:rls}{facteur d'oubli} transforme cette demi-vie
en mémoire de l'estimateur qui négocie le signal.
\item L'heure, le détecteur et son seuil, la trajectoire de la statistique, la version du modèle, et ce qui a été fait~: réentraînement,
dérogation, ou rejet, avec le motif.
\item Comme un changement dans la distribution des variables, les gains restant inchangés~; les détecteurs sur les variables (chapitre~27) le voient,
les détecteurs sur les gains ne le voient pas avant que cela ne coûte.
\item Aussi vite que la vérité se déplace et pas davantage, ce qu'on mesure sur l'historique et ce qu'on vérifie sur les gains.
\end{enumerate}
\end{solution}

\begin{solution}{iq:ml:online-learning-and-drift:1}
La \omterm{def:ml:online-learning-and-drift:drift}{dérive conceptuelle} (concept drift) modifie $\P(y\mid x)$~: un signal est évincé par l'encombrement (crowding), ou une modification des frais inverse le signe de l'effet
d'une variable de flux d'ordres. Le \omterm{def:ml:online-learning-and-drift:drift}{décalage des covariables} (covariate shift) ne modifie que $\P(x)$~: un régime de volatilité déplace les variables dans des plages que le modèle a rarement vues.
\iqlookfor{les deux distributions et un exemple correct de chacune.}
\end{solution}

\begin{solution}{iq:ml:online-learning-and-drift:2}
Minimiser $\sum_s\lambda^{t-s}(y_s - x_s^\top b)^2$~; les équations normales $R_tb = r_t$ avec $R_t = \lambda R_{t-1} + x_tx_t^\top$
et $r_t = \lambda r_{t-1} + x_ty_t$~; la formule de Sherman--Morrison appliquée à $P_t = R_t^{-1}$ donne le gain, la mise à jour des coefficients et la
mise à jour de la covariance de la \cref{def:ml:online-learning-and-drift:rls}. $\lambda$ fixe la mémoire, soit environ $1/(1-\lambda)$
observations.
\iqlookfor{le critère pondéré, la mise à jour de rang un et l'interprétation de la mémoire.}
\end{solution}

\begin{solution}{iq:ml:online-learning-and-drift:3}
Sur des données sans changement, on choisit le seuil qui donne un taux de fausses alarmes sur lequel le desk peut agir~; on injecte ensuite les changements qui
comptent et on mesure le délai~; on rapporte les deux. Il faut vérifier sur un historique réservé que le taux tient.
\iqlookfor{calibration sur des données sans changement, délai mesuré sur des changements non vus, et le compromis entre les deux.}
\end{solution}

\begin{solution}{iq:ml:online-learning-and-drift:4}
Évaluer chaque \omterm{def:ml:online-learning-and-drift:schedule}{calendrier de réentraînement} en prequential sur l'historique, en comptant les coûts de réajustement~; estimer la vitesse de déplacement de la relation (par balayage du
\omterm{def:ml:online-learning-and-drift:rls}{facteur d'oubli} ou de la largeur de fenêtre)~; mettre en balance le gain et le risque opérationnel lié à des déploiements plus fréquents~; garder le
calendrier comme plancher et ajouter un détecteur.
\iqlookfor{une comparaison prequential, les coûts et une estimation de la dérive.}
\end{solution}

\begin{solution}{iq:ml:online-learning-and-drift:5}
Il cumule les écarts de la valeur sous sa moyenne courante, diminués d'une tolérance, et déclenche une alarme lorsque la somme remonte
d'un seuil au-dessus de son minimum courant. Faiblesses~: la moyenne courante porte sur tout l'historique depuis la dernière remise à zéro, si bien qu'après une
longue période calme, elle réagit lentement~; elle ne détecte bien qu'une seule direction et une seule ampleur de changement, celles pour lesquelles sa tolérance et
son seuil ont été réglés.
\iqlookfor{la statistique et une vraie faiblesse (mémoire de la moyenne courante, ou réglage sur un seul changement).}
\end{solution}

\begin{solution}{iq:ml:online-learning-and-drift:6}
Voir la \cref{prop:ml:online:kalman}~: en multipliant la covariance d'état par $1/\lambda$ à l'étape de prédiction, on retrouve la méthode RLS. Ainsi
$\lambda$ traduit la vitesse à laquelle les coefficients bougent par rapport au bruit~; on le choisit en l'évaluant sur l'historique, ou on estime le
bruit d'état par maximum de vraisemblance, puis on en déduit le $\lambda$ implicite.
\iqlookfor{l'étape de prédiction mise en correspondance avec l'oubli, et la conséquence pour le choix de $\lambda$.}
\end{solution}
